% Page budget: 1.55 pages | Owns: augmentation topology, additional states, encoder, and closed-loop objective.
% WRITING GUIDE
% Purpose: Define the final augmentation architecture and explain how it is trained inside the known feedback loop.
% Start here: Current implementation decisions D-129, D-140, D-141, D-180, D-220, D-233, and D-237 in docs/decisions.md.
% Supporting material: Quinten's 18 September outline feedback, the Meeting-audit synthesis, Thesis-writeup/Documentation/TRAINING-DESIGN.md, and docs/writeup figures.
% Mandatory papers: Read Hoekstra's 2025 LFR augmentation paper, the 2026 first-principles augmentation paper, the encoder-initialization paper, and Kessels' Chapter 5 before writing this section.
% Research-plan anchor: Preserve Aspect 2's dynamic parallel augmentation objective; document the departure from open-loop training and earlier state routing.
% Current decisions: Separate physical and additional states, keep the controller outside the model, score the complete training window without burn-in, and use one network for the physical correction and added-state update.
% Choices to justify: Final reported routing, added-state dimension, ANN width and depth, zero initialization, encoder history, closed-loop objective, rollout length, full-window scoring, and validation horizon. Use the configuration of the reported thesis run, not a superseded routing experiment.
% Horizon check: Treat encoder history, scored training window, joint-estimation horizon, and free-run validation horizon separately; relate candidates to relevant stable modes and test sensitivity to slower dynamics and free-integrator accumulation.
% Implementation audit: Read scripts/gantry/gantry_interconnect_dynamic.py, gantry_dynamic/{model,config,controller,data,training}.py, and model_augmentation/fit_systems/{interconnect,closed_loop,pre_encoder}.py.
% Reference check: Verify sources for parallel LFR augmentation, encoder initialization, closed-loop state extension, and any claimed architectural precedent against the original paper or thesis.
% Cautions: Earlier routing plans are superseded; distinguish the truth-only hidden absorber from learned states; closed-loop fit alone does not prove autonomous plant recovery.
% Figures: Absorber schematic, author's updated architecture, and thesis-owned common-reference loops. No residual panel or horizon figure.
% Section is finished when: Every signal has a defined frame and timing, the RK4 boundary is explicit, and the described architecture matches the final code.
%
% SOURCE MAP
% Hoekstra et al. 2025 EJC: dynamic-parallel topology, additional states, encoder-based truncated simulation.
% Hoekstra et al. 2026 Automatica submission: extended LFR formulation, baseline-preserving model/encoder initialization, joint identification context.
% Hoekstra et al. 2026 encoder paper: reconstructability-based W^b initialization; it does not identify W^a or test dynamic augmentation.
% Kessels 2025 thesis: direct closed-loop rollout, controller equations, and reconstruction of the machine controller state for each window.
% Drenth 2025 thesis: LPV-specific augmentation precedent; not the unrestricted MLP structure used here.
% This project: gantry specialization, residual-loop implementation, exact routing, RK4 boundary, and verification.
\section{Dynamic LPV-LFR Augmentation}\label{sec:augmentation}

The baseline of Section~\ref{sec:lfr} describes the rigid-body dynamics of
the gantry, but not every effect of the machine.
Garc\'ia-Herreros et al.\ neglect the resonance of the supporting base, the
vibration of the cross-arm on its flexible plates, the detent forces of the
linear motors, and the variation of friction along the
stroke~\cite[Sec.~2.4]{garcia2013model}, and our baseline \eqref{eq:eom}
also leaves out the Coulomb friction that their model contains.
We therefore augment the baseline with a learned component and estimate the
physical parameters $\theta_{\mathrm{base}}$, the network parameters
$\theta_{\mathrm{aug}}$ and the encoder parameters $\eta$ jointly from
records measured under feedback.
To this end, we specify the augmented model
(Section~\ref{sec:aug_structure}), the encoder that supplies the initial
state of each training window (Section~\ref{sec:aug_encoder}), the
simulation of a window through the known feedback controller
(Section~\ref{sec:aug_closed_loop}), and the identification criterion
(Section~\ref{sec:aug_objective}).
Each record contains the actuator forces $\uact$ and the measured positions
$\qact$, sampled with period $T_s$, and we write $u_{\mathrm{act},k}$ and
$y_k$ for their values at sample $k$.

\subsection{Augmented model}\label{sec:aug_structure}
An omitted vibration mode requires additional model order.
Refitting $\theta_{\mathrm{base}}$ cannot add this order, and neither can a
static correction of the baseline's state update, because both only change
the dynamics of the existing states.
A dynamic augmentation adds states for the missing
dynamics~\cite[Sec.~2]{hoekstra2025lfr}, and Drenth formulates such an
augmentation for a self-scheduled LPV-LFR baseline, with the augmented part
itself an LPV-LFR~\cite[Sec.~5.2]{drenth2025thesis}.
We use the dynamic-parallel structure of~\cite[Table~1]{hoekstra2026lfr},
with a neural network as the learned part, which gives the augmented model
\begin{equation}
\begin{aligned}
 \tilde x_{k+1}
 &=f_{\mathrm{base}}(\theta_{\mathrm{base}},\tilde x_k,\hat u_{\mathrm{act},k})
   +f_{\mathrm{aug}}(\theta_{\mathrm{aug}},\hat x_k,\hat u_{\mathrm{act},k}),\\
 \bar x_{k+1}
 &=g_{\mathrm{aug}}(\theta_{\mathrm{aug}},\hat x_k,\hat u_{\mathrm{act},k}),\\
 \hat y_k&=h_{\mathrm{base}}(\tilde x_k),
\end{aligned}
\label{eq:aug_transition}
\end{equation}
where $\tilde x\in\R^6$ is the physical state of Section~\ref{sec:lfr},
$\bar x\in\R^{n_{x_a}}$ collects the additional states,
$\hat x=\col(\tilde x,\bar x)$, $h_{\mathrm{base}}$ is the output map
defined in \eqref{eq:aug_norm}, and $\hat u_{\mathrm{act},k}$ is the actuator
force applied to the model during the closed-loop simulation of
Section~\ref{sec:aug_closed_loop}.
The baseline transition $f_{\mathrm{base}}$ keeps its physical parameters,
and the learned functions $f_{\mathrm{aug}}$ and $g_{\mathrm{aug}}$ act in
parallel to it (Fig.~\ref{fig:aug_structure}).

\begin{figure}[t]
 \centering
 \resizebox{\columnwidth}{!}{\input{figures/tex/augmentation_structure}}
 \caption{Augmented model \eqref{eq:aug_transition}. The physics step
 $f_{\mathrm{base}}$ and the network $\phi_{\mathrm{aug}}$ act in parallel
 on $\hat x_k$; $\phi_{\mathrm{aug}}$ supplies $f_{\mathrm{aug}}$ and
 $g_{\mathrm{aug}}$, and the output map is not learned
 ($h_{\mathrm{aug}}=0$). The encoder $\psi$ sets $\hat x$ at each window
 start. The input $u_k$ is the actuator force; during training it is the
 closed-loop input $\hat u_{\mathrm{act},k}$ of \eqref{eq:aug_controller}.}
 \label{fig:aug_structure}
\end{figure}

The baseline transition $f_{\mathrm{base}}$ is one step of the classical
fourth-order Runge-Kutta (RK4) method applied to the state equation of
\eqref{eq:eom}, written as
$\dot{\tilde x}=F(\theta_{\mathrm{base}},\tilde x,\uact)$.
The actuator force is held constant over the step.
Each of the four stages evaluates $F$ at its own state through the LFR of
Section~\ref{sec:lfr}, so the scheduling variable $Y$ moves within the step
and $f_{\mathrm{base}}$ remains self-scheduled.
The physics is thus integrated in continuous time, whereas
$f_{\mathrm{aug}}$ is added to the result of the step and is a discrete-time
correction.

The inputs and states of the gantry differ in scale by orders of magnitude,
and normalising the data to zero mean and unit standard deviation improves
the estimation of models that contain neural
networks~\cite[Sec.~5.3]{hoekstra2026lfr}.
We therefore normalise every signal with the sample mean and standard
deviation of the training records, and evaluate the baseline in physical
coordinates,
\begin{equation}
\begin{aligned}
 \tilde x_{\mathrm n}&=S_x^{-1}(\tilde x-\mu_x),\quad
 u_{\mathrm n}=S_u^{-1}(\uact-\mu_u),\quad
 y_{\mathrm n}=S_y^{-1}(y-\mu_y),\\
 f_{\mathrm{base},\mathrm n}(\theta_{\mathrm{base}},\tilde x_{\mathrm n},u_{\mathrm n})
 &=S_x^{-1}\bigl(f_{\mathrm{base}}(\theta_{\mathrm{base}},
   \mu_x+S_x\tilde x_{\mathrm n},\mu_u+S_uu_{\mathrm n})-\mu_x\bigr),\\
 h_{\mathrm{base}}(\tilde x_{\mathrm n})&=C_{\mathrm n}\tilde x_{\mathrm n},\qquad
 C_{\mathrm n}=S_y^{-1}\begin{bmatrix}P^\top & 0\end{bmatrix}S_x,
\end{aligned}
\label{eq:aug_norm}
\end{equation}
where $\mu_x,\mu_u,\mu_y$ are the sample means and $S_x,S_u,S_y$ the diagonal
matrices of sample standard deviations.
The state statistics are computed from the positions $P^{-\top}y_k$ and their
first differences, so they come from measured signals rather than from a
simulation of the baseline as in~\cite[Sec.~5.3]{hoekstra2026lfr}.
Because $\mu_y=\begin{bmatrix}P^\top & 0\end{bmatrix}\mu_x$, the output map
$C_{\mathrm n}$, the kinematic map of \eqref{eq:Ptransform} in normalised
coordinates, carries no offset.
The transformation changes the coordinates of $f_{\mathrm{base}}$ but not its
dynamics.
The network, the encoder and the objective all operate in these coordinates,
and from here on we drop the subscript $\mathrm n$.

The two learned functions are the outputs of one multilayer perceptron (MLP)
with tanh hidden layers,
\begin{equation}
 \begin{bmatrix}
  f_{\mathrm{aug}}(\theta_{\mathrm{aug}},\hat x_k,\hat u_{\mathrm{act},k})\\
  g_{\mathrm{aug}}(\theta_{\mathrm{aug}},\hat x_k,\hat u_{\mathrm{act},k})
 \end{bmatrix}
 =\phi_{\mathrm{aug}}(\theta_{\mathrm{aug}},z_k),\qquad
 z_k=\col(\hat x_k,\hat u_{\mathrm{act},k}),
 \label{eq:aug_mlp}
\end{equation}
where $\theta_{\mathrm{aug}}$ collects the weights and biases of the MLP.
The first six outputs form $f_{\mathrm{aug}}$ and correct every
physical-state row, the positions and velocities of $X$ and $Y$ included;
the last $n_{x_a}$ outputs form $g_{\mathrm{aug}}$ and give the next
additional state.
The rows of $X$ and $Y$ are included because the omitted dynamics couple
into the $X$ and $Y$ responses, which a correction restricted to the yaw rows
cannot represent, and the learning components
of~\cite[Sec.~2]{hoekstra2025lfr} likewise act on the entire model state.
As a consequence, a correction written into the rows of $X$ and $Y$, which
have no stiffness in \eqref{eq:damping_stiffness_matrices}, is integrated
without a restoring force; Section~\ref{sec:aug_closed_loop} returns to
this.

The output map is not augmented, so the measured positions are read from the
physical state alone, and the additional states reach the output only
through $f_{\mathrm{aug}}$, one step after they are formed.
Within one step, no learned function feeds the input of another, so the
interconnection graph of \eqref{eq:aug_transition} is acyclic.
$F$ is smooth wherever $M(Y)$ is invertible, which the parameterisation of
Section~\ref{sec:params} guarantees over the whole stroke, and the tanh
activations are smooth as well.
Both components are therefore twice differentiable, and the augmented model
is well posed by~\cite[Thm.~9]{hoekstra2026lfr}.

Training starts from the baseline.
The output layer of the MLP is initialised at zero,
\begin{equation}
 W_L=0,\quad b_L=0
 \quad\Longrightarrow\quad
 f_{\mathrm{aug}}=0,\quad g_{\mathrm{aug}}=0 ,
 \label{eq:aug_zero_init}
\end{equation}
where $W_L$ and $b_L$ are the weight matrix and bias of that layer.
The augmented model then reproduces the baseline from the same physical
initial state, and $\bar x_{k+1}=0$ at every step; this is the baseline
behaviour at initialisation of~\cite[Eq.~(29)]{hoekstra2026lfr} with
$g_{\mathrm{aug}}=0$.
From this start, $f_{\mathrm{aug}}$ can also learn changes that a different
$\theta_{\mathrm{base}}$ would produce, so the split between the two
components is not unique~\cite[Sec.~5.2]{hoekstra2026lfr}, and
Section~\ref{sec:obc} constrains it.
The number of additional states and the size of the MLP are given in
Section~\ref{sec:setup}.

\subsection{Encoder and initial state}\label{sec:aug_encoder}
Each training window simulates \eqref{eq:aug_transition} from an initial
state, but only the positions are measured; the velocities and the
additional states are not.
Following Beintema et al.~\cite{beintema2023subnet}, an encoder estimates
this state from the recent input and output history,
\begin{equation}
 \hat x_{\tau|\tau}
 =\psi_\eta(\xi_\tau),\qquad
 \xi_\tau=\col\bigl(y_{\tau-n:\tau},\,u_{\mathrm{act},\tau-n:\tau}\bigr),
 \label{eq:aug_encoder}
\end{equation}
where $\tau$ is the window start, $n$ the length of the history, and
$y_{\tau-n:\tau}$ stacks the samples $y_{\tau-n},\dots,y_\tau$, and likewise
for the input.
The history includes the current sample, as the reconstructability map
of~\cite[Eq.~(15)]{hoekstra2026encoder} does.
The encoder estimates all states, the measured positions included.

The encoder is a linear map with a nonlinear
correction~\cite[Eq.~(8)]{hoekstra2026encoder},
\begin{equation}
 \psi_\eta(\xi_\tau)
 =\begin{bmatrix}W^b\\ W^a\end{bmatrix}\xi_\tau+\tilde\psi(\xi_\tau),
 \label{eq:aug_enc_lin}
\end{equation}
where $W^b$ and $W^a$ give the physical and the additional states, and
$\tilde\psi$ is an MLP; $\eta$ collects $W^b$, $W^a$ and the parameters of
$\tilde\psi$.
We initialise the rows $W^b$ from the baseline.
Linearising \eqref{eq:eom} about rest at $Y=0$, with the nominal parameters
of Appendix~\ref{app:params}, gives
\begin{equation}
 A_c=\begin{bmatrix}0 & I\\ -M(0)^{-1}K & -M(0)^{-1}C\end{bmatrix},\quad
 B_c=\begin{bmatrix}0\\ M(0)^{-1}P\end{bmatrix},\quad
 C_d=\begin{bmatrix}P^\top & 0\end{bmatrix},
 \label{eq:aug_lin}
\end{equation}
and $(A_d,B_d)$ is the zero-order-hold discretisation of $(A_c,B_c)$ over
$T_s$, consistent with the held actuator force.
All matrices are normalised as the signals in \eqref{eq:aug_norm}.

For this linear model, $W^b$ is the reconstructability
map~\cite[Eqs.~(16), (17)]{hoekstra2026encoder}
\begin{equation}
 W^b=\begin{bmatrix}
  A_d^nO_n^{+} & \;R_n-A_d^nO_n^{+}T_n
 \end{bmatrix},
 \label{eq:aug_wb}
\end{equation}
where $O_n$ is the observability matrix of $(A_d,C_d)$ over the $n+1$
samples of the history, $T_n$ the Toeplitz matrix of its input-to-output
responses, $R_n$ the input-to-state map over the same samples, and the two
column blocks act on the output and input parts of $\xi_\tau$.
The measured positions make $(A_d,C_d)$ observable, so $O_n$ has full column
rank and $O_n^{+}$ is a left inverse.
The linearisation only sets the starting value of $W^b$, because the encoder
is trained jointly with the model.
The rows $W^a$ have no baseline counterpart and are drawn randomly, and the
output layer of $\tilde\psi$ starts at zero.
By \eqref{eq:aug_zero_init}, the encoded additional states then do not
affect the output at initialisation; they act on it only once the output
layer of $\phi_{\mathrm{aug}}$ departs from zero.

\subsection{Closed-loop rollout}\label{sec:aug_closed_loop}
The records are measured under feedback, and the model is meant to predict
the machine under feedback: its response to a given reference, feedforward
and controller, including a controller other than the one in the data.
The criterion of~\cite[Eq.~(22)]{hoekstra2026lfr} instead simulates the
model in open loop, driven by the recorded input.
Kessels trains an augmented model with the known feedback controller inside
each truncated simulation and propagates gradients through
it~\cite[Eqs.~(5.12), (5.13c), (5.13d)]{kessels2025ai}.
We adopt this closed-loop rollout and keep the controller outside the
learned model, so that it can be replaced in evaluation, as Kessels does
with a retuned controller~\cite[Sec.~5.3.2.3]{kessels2025ai}.

The baseline gives a second reason.
$X$ and $Y$ have no stiffness in \eqref{eq:damping_stiffness_matrices}, so
the baseline, frozen at any $Y$, has two eigenvalues at the origin.
In an open-loop rollout, a position error caused by a velocity error
therefore persists, and a correction of $f_{\mathrm{aug}}$ on these rows
accumulates in the same way.
Inside the loop, the controller acts on such errors as it does on the
machine, provided the model in feedback with the controller is stable.
The same feedback also attenuates model errors at the output, so a good
closed-loop fit alone does not show that the plant model is correct.

The recorded loop is driven by a reference $r$ and a feedforward
$u^{\mathrm{ff}}$ through a linear time-invariant controller with state $s$
and matrices $(A_{\mathcal K},B_{\mathcal K},C_{\mathcal K},D_{\mathcal K})$.
The model loop receives the same $r$, $u^{\mathrm{ff}}$ and controller, and
differs only in its output $\hat y$ (Fig.~\ref{fig:aug_closed_loop}):
\begin{equation}
\begin{aligned}
 s_{k+1}&=A_{\mathcal K}s_k+B_{\mathcal K}(r_k-y_k), &
 u_{\mathrm{act},k}&=C_{\mathcal K}s_k+D_{\mathcal K}(r_k-y_k)+u^{\mathrm{ff}}_k,\\
 \hat s_{k+1}&=A_{\mathcal K}\hat s_k+B_{\mathcal K}(r_k-\hat y_k), &
 \hat u_{\mathrm{act},k}&=C_{\mathcal K}\hat s_k+D_{\mathcal K}(r_k-\hat y_k)+u^{\mathrm{ff}}_k ,
\end{aligned}
\label{eq:aug_direct_controller}
\end{equation}
where both loops feed back the error of the output with respect to the
reference, and $D_{\mathcal K}\neq0$, so the controller has direct
feedthrough.
Subtracting the recorded loop from the model loop removes $r$ and
$u^{\mathrm{ff}}$,
\begin{equation}
\begin{aligned}
 x^c_{k+1}&=A_{\mathcal K}x^c_k+B_{\mathcal K}(y_k-\hat y_k),\qquad x^c_\tau=0,\\
 \hat u_{\mathrm{act},k}&=u_{\mathrm{act},k}+C_{\mathcal K}x^c_k+D_{\mathcal K}(y_k-\hat y_k),
\end{aligned}
\label{eq:aug_controller}
\end{equation}
where $x^c=\hat s-s$ is the difference of the controller states.
The model thus receives the recorded input plus the controller's response to
its own output error, and needs neither the reference nor the feedforward.
This form is exact when both loops apply the same linear controller, not
scheduled on the model state, with the same feedforward and without
saturation.

\begin{figure}[t]
 \centering
 \includegraphics[width=\columnwidth]{closed_loop_same_reference.pdf}
 \caption{Recorded plant loop (top) and augmented-model loop (bottom),
 driven by one shared reference $r_k$ and feedforward $u_k^{\mathrm{ff}}$,
 with the same feedback controller. The model output is evaluated
 before the input advances its state to the next sample.}
 \label{fig:aug_closed_loop}
\end{figure}

The condition $x^c_\tau=0$ starts the model's controller in the state of the
recorded controller, $\hat s_\tau=s_\tau$, without reconstructing that state,
because the recorded input already contains its effect.
Kessels instead reconstructs the controller state from the measured output,
the reference and the controller~\cite[Remark~5.4]{kessels2025ai}.
Model error from before $\tau$ is therefore not carried into the window.
Because the model has no feedthrough and the controller does, each sample is
evaluated in a fixed order: $\hat y_k$ from the state, then
$\hat u_{\mathrm{act},k}$ from \eqref{eq:aug_controller}, and then both
states advance.
This order closes the loop without an algebraic loop and without a delay of
one sample.

\subsection{Identification criterion}\label{sec:aug_objective}
A simulation over a whole record is one long sequential computation.
Following~\cite[Sec.~5.1]{hoekstra2026lfr}, we therefore split the records
into short overlapping windows, which allows batch optimisation and
increases data efficiency.
Over a set of window starts $\mathcal T$ and a window length $n_f$, the
parameters $\vartheta=(\theta_{\mathrm{base}},\theta_{\mathrm{aug}},\eta)$
minimise the output error
\begin{equation}
 V(\vartheta)=\frac{1}{|\mathcal T|\,n_f}
 \sum_{\tau\in\mathcal T}\sum_{\ell=0}^{n_f-1}
 \norm{y_{\tau+\ell}-\hat y_{\tau+\ell|\tau}}_2^2
 \label{eq:aug_loss}
\end{equation}
subject to \eqref{eq:aug_transition}, \eqref{eq:aug_encoder} and
\eqref{eq:aug_controller}, where $\hat y_{\tau+\ell|\tau}$ is the output
$\ell$ samples after the window start, simulated from
$\hat x_{\tau|\tau}$.
This is the truncated objective of~\cite[Eq.~(22)]{hoekstra2026lfr} with
each window simulated in closed loop, as in~\cite[Eq.~(5.12)]{kessels2025ai},
and with the physical parameters, the network and the encoder estimated
jointly while the controller is fixed.
Unlike~\cite[Eq.~(27)]{hoekstra2026lfr}, $V$ contains no term that
penalises the deviation of $\theta_{\mathrm{base}}$ from nominal values.
Such a term bounds the drift of the parameters but does not remove the
non-uniqueness of Section~\ref{sec:aug_structure}, which
Section~\ref{sec:obc} addresses instead.

Every sample of the window is scored, from $\ell=0$ on, as in both cited
objectives.
The error of the encoded initial state therefore remains part of $V$, and
the encoder is trained on it.
The objective involves three horizons that must be kept apart: the encoder
history $n$, the window length $n_f$, and the horizon of the free-run
simulation used for validation.
The window must span the time over which a model error remains visible in
the closed-loop output, because a shorter window scores only part of its
effect; Section~\ref{sec:setup} gives the three values.
